\originalchapter{8.701 作业一:粒子与对称性}
题源: Fall 2020 Problem Set 1, 2020年9月10日发布. 五题均有对应官方答案. 本章数值表采用官方答案所用的近似值, 是课程数据版本, 不是最新 PDG 汇总.

\begin{exercise}\label{ex:8701-1-1}
原题 1. 查阅 $\pi^\pm,\pi^0,K^\pm,K_S^0,K_L^0,\mu^\pm,e^\pm$, 质子及中子的质量、平均寿命和最大分支的主要衰变模式. 假定速度为 $c$ 且忽略时间膨胀, 各粒子平均走多远?
\end{exercise}
\begin{solution}
官方解答重编. 题设的 $c\tau$ 是固有衰变长度的比较指标, 有质量粒子不能真的以 $c$ 运动; 实际平均实验室距离为 $\beta\gamma c\tau$.
\begin{center}\begin{tabular}{lllll}\toprule
粒子&质量(MeV/$c^2$)&$\tau$(s)&主要模式&$c\tau$(m)\\\midrule
$\pi^\pm$&$140$&$2.6\times10^{-8}$&$\mu^\pm\nu$&$7.8$\\
$\pi^0$&$135$&$8.5\times10^{-17}$&$\gamma\gamma$&$2.6\times10^{-8}$\\
$K^\pm$&$494$&$1.2\times10^{-8}$&$\mu^\pm\nu$&$3.6$\\
$K_S^0$&$498$&$9\times10^{-11}$&$\pi^+\pi^-$&$0.027$\\
$K_L^0$&$498$&$5\times10^{-8}$&$\pi e\nu$ (两电荷支合计)&$15$\\
$\mu^\pm$&$106$&$2.2\times10^{-6}$&$e\nu\bar\nu$&$660$\\
$e^\pm$&$0.511$&稳定&无自发衰变&无有限值\\
$p$&$938$&本课程视为稳定&未观察自发衰变&无有限值\\
$n$&$940$&$880$&$pe^-\bar\nu_e$&$2.6\times10^{11}$\\\bottomrule
\end{tabular}\end{center}
中子的条目指自由中子. 正电子在物质中能湮灭, 这与真空中的固有不稳定性不同.
\end{solution}

\begin{exercise}\label{ex:8701-1-2}
原题 2. 静止带电 pion 可以衰变到电子或 muon 及相应中微子. (a) 画出末态动量与自旋, 解释最大宇称破坏和哪条分支较大. (b) 取中微子无质量, 证明 $p_\ell=(m_\pi^2-m_\ell^2)/(2m_\pi)$, $E_\ell=(m_\pi^2+m_\ell^2)/(2m_\pi)$. 给定 $m_\pi=139.57$, $m_\mu=105.66$, $m_e=0.511$, 单位 MeV.
\end{exercise}
\begin{solution}
官方解答重编. 以 $\pi^-$ 为例, 让 $\ell^-$ 向左、$\bar\nu_\ell$ 向右. 反中微子右螺旋度使其自旋向右; 自旋零母态要求轻子自旋向左. 两个自旋反向, 轻子的正螺旋度需要质量引起的左手征成分. P 反转动量却不反转自旋, 会把反中微子换成弱作用不产生的螺旋度, 所以分布不具 P 对称.
\begin{center}\begin{tikzpicture}
\draw[->,thick](-0.4,0)--(-2,0)node[left]{$\mathbf p_\ell$};\draw[->,thick](0.4,0)--(2,0)node[right]{$\mathbf p_{\bar\nu}$};
\draw[->,structurecolor](-0.4,-0.5)--(-1.5,-0.5)node[left]{$\mathbf s_\ell$};\draw[->,structurecolor](0.4,-0.5)--(1.5,-0.5)node[right]{$\mathbf s_{\bar\nu}$};\node at(0,0.4){$\pi^-$ 静止};
\end{tikzpicture}\end{center}
宽度比由第 \ref{ch:spinor} 章为约 $1.3\times10^{-4}$, 所以 muon 分支大. 两体能量守恒 $m_\pi=p+\sqrt{p^2+m_\ell^2}$, 平方一次得到所求 $p$, 再由 $E_\ell=m_\pi-p$ 得能量. 数值 muon 动量约 $29.8\,\mathrm{MeV}$、总能量 $109.8\,\mathrm{MeV}$; 电子动量和总能量均约 $69.8\,\mathrm{MeV}$.
\end{solution}

\begin{exercise}\label{ex:8701-1-3}
原题 3. Higgs 对费米子的耦合正比质量, 部分宽度正比耦合平方. 简化为只能衰变到 bottom、charm、tau 和 muon, 求各分支比.
\end{exercise}
\begin{solution}
官方解答重编. 取源答案 $m_b=4.18$, $m_c=1.28$, $m_\tau=1.78$, $m_\mu=0.11\,\mathrm{GeV}$. 权重为 $3m_b^2,3m_c^2,m_\tau^2,m_\mu^2$, 分母总和约 $60.51\,\mathrm{GeV}^2$. 依次为 $86.62\%,8.12\%,5.24\%,0.0200\%$. 颜色因子必须保留; 这些是题设模型的相对比例, 不是完整 Higgs 实际分支比.
\end{solution}

\begin{exercise}\label{ex:8701-1-4}
原题 4. 分别讨论 (a) 位置 $\mathbf r$, (b) 动量 $\mathbf p$, (c) 轨道角动量 $\mathbf L$, (d) 自旋 $\mathbf s$, (e) 静电场 $\mathbf E$, (f) 静磁场 $\mathbf B$, (g) $\mathbf s\cdot\mathbf E$, (h) $\mathbf s\cdot\mathbf B$ 在 P 与 T 下的变换.
\end{exercise}
\begin{solution}
官方解答重编. 各量的 P、T 奇偶性依次为
$( -,+),(-,-),(+,-),(+,-),(-,+),(+,-),(-,-),(+,+)$.
例如 $\mathbf L=\mathbf r\times\mathbf p$ 在 P 下两个负号相消, T 下只有动量变号; 磁场由运动电荷产生, 与角动量同类. 电偶极项在 P、T 下都变号, 磁偶极项都不变. 判断包含外场的能量时应同时变换系统和外场.
\end{solution}

\begin{exercise}\label{ex:8701-1-5}
原题 5. 电磁、弱、强三类相互作用各保持哪些量: (a) 轻子数, (b) 强同位旋 $I_3$, (c) 奇异数 S, (d) 宇称, (e) CP?
\end{exercise}
\begin{solution}
官方入门表重编及条件补充. 按通常微扰标准模型、忽略强 CP、取总轻子数而非逐味轻子数: 三种作用都保持总轻子数; $I_3$, S, P, CP 在强和电磁作用中保持, 在弱作用中一般不保持. 电磁作用保持 $I_3$ 并不代表保持完整 SU(2) 同位旋. 中微子混合使各味轻子数变化; Majorana 质量或其他扩展可破坏总轻子数, 强 $\theta$ 项也会改变 CP 判断. 因而表中结论有声明的适用层次.
\end{solution}
